\section{Modellierung eines Datenmodells für Datenbanksysteme}

\subsection*{Ziel}
Ein \textbf{Datenmodell} beschreibt die Struktur, Beziehungen und Regeln der Daten, 
die in einem IT-System gespeichert und verarbeitet werden. 
Das Ziel besteht darin, eine \textbf{konsistente, redundanzfreie und performante Abbildung der Realwelt} in Datenbankform zu schaffen.

\subsection*{1. Phasen der Datenmodellierung}

\subsubsection*{1.1 Konzeptionelles Datenmodell (Fachliches Modell)}
\begin{itemize}
	\item \textbf{Ziel:} Abbilden der realen Welt aus Sicht des Fachbereichs.
	\item \textbf{Werkzeug:} Entity-Relationship-Diagramm (ERD) oder UML-Klassendiagramm.
	\item \textbf{Elemente:}
	\begin{itemize}
		\item \textbf{Entitäten} (z.\,B. \textit{Kunde}, \textit{Produkt}, \textit{Bestellung})
		\item \textbf{Attribute} (z.\,B. \textit{Kundenname}, \textit{Preis})
		\item \textbf{Beziehungen} (z.\,B. \textit{Kunde gibt Bestellung auf})
	\end{itemize}
\end{itemize}
\textbf{Ergebnis:} Ein fachliches, technologieunabhängiges ER-Diagramm.

\subsubsection*{1.2 Logisches Datenmodell}
\begin{itemize}
	\item \textbf{Ziel:} Präzisierung des konzeptionellen Modells auf Datenbankebene.
	\item \textbf{Werkzeug:} Logisches ER-Diagramm (z.\,B. Chen- oder Crow’s-Foot-Notation).
	\item \textbf{Schritte:}
	\begin{itemize}
		\item Definition von Primär- und Fremdschlüsseln.
		\item Festlegung der Beziehungstypen (1:1, 1:n, m:n).
		\item Normalisierung (mindestens bis zur 3. Normalform, oft 3NF oder BCNF).
	\end{itemize}
\end{itemize}
\textbf{Ergebnis:} Ein strukturierter, logischer Entwurf, der sich direkt in Tabellen übersetzen lässt.

\subsubsection*{1.3 Physisches Datenmodell}
\begin{itemize}
	\item \textbf{Ziel:} Umsetzung in einem konkreten Datenbanksystem.
	\item \textbf{Werkzeug:} SQL-DDL-Skripte (\texttt{CREATE TABLE ...}), Datenbank-Design-Tools.
	\item \textbf{Schritte:}
	\begin{itemize}
		\item Definition von Datentypen (z.\,B. \texttt{VARCHAR(255)}, \texttt{INT}, \texttt{DATE})
		\item Indizes, Constraints, Trigger
		\item Performance-Überlegungen (Partitionierung, Denormalisierung, Materialized Views)
		\item Physische Speicheroptionen (z.\,B. Tablespaces, Sharding)
	\end{itemize}
\end{itemize}
\textbf{Ergebnis:} Ein lauffähiges Datenbankschema, bereit zur Implementierung.

\subsection*{2. Vorgehensweise in der Praxis}
\begin{enumerate}[label=\arabic*.]
	\item \textbf{Anforderungsanalyse:} Erheben, welche Informationen das System speichern und verarbeiten muss.
	\item \textbf{Erstellen des konzeptionellen Modells:} Entitäten, Attribute und Beziehungen skizzieren.
	\item \textbf{Validierung mit Fachanwendern:} Prüfen, ob das Modell die Realität korrekt abbildet.
	\item \textbf{Überführung ins logische Modell:} Schlüsselfelder, Kardinalitäten und Normalisierung definieren.
	\item \textbf{Erstellung des physischen Modells:} Umsetzung für ein konkretes DBMS (z.\,B. PostgreSQL).
	\item \textbf{Optimierung und Pflege:} Performance-Tuning, Versionierung, Änderungsmanagement.
\end{enumerate}

\subsection*{3. Tools zur Datenmodellierung}
\begin{tabular}{@{}p{4cm}p{10cm}@{}}
	\toprule
	\textbf{Typ} & \textbf{Beispiele} \\ \midrule
	Professionelle Werkzeuge & ER/Studio, SAP PowerDesigner, Oracle SQL Developer Data Modeler \\
	Open Source / Kostenlos & dbdiagram.io, draw.io, Lucidchart, DBeaver, pgModeler \\
	Codebasiert & Prisma (Node.js), SQLAlchemy (Python), Hibernate (Java ORM) \\
	\bottomrule
\end{tabular}

\subsection*{4. Beispiel eines logischen Modells}
Ein einfaches logisches Modell für einen Online-Shop:
\begin{verbatim}
	Kunde (KundenID PK, Name, E-Mail)
	Bestellung (BestellID PK, Datum, KundenID FK)
	Produkt (ProduktID PK, Bezeichnung, Preis)
	Bestellposition (BestellID FK, ProduktID FK, Menge)
\end{verbatim}

\noindent Beziehungen:
\begin{itemize}
	\item Kunde — Bestellung: 1:n
	\item Bestellung — Bestellposition: 1:n
	\item Produkt — Bestellposition: 1:n
\end{itemize}